La palta Hass se ha consolidado como un producto agrícola de alto valor comercial por el crecimiento de su consumo en Norteamérica, Europa y Asia. Sin embargo, su precio de exportación presenta variaciones asociadas con la estacionalidad de las cosechas, la oferta de los países competidores, el destino comercial, el volumen disponible y los costos logísticos. Estas fluctuaciones afectan la estabilidad económica de productores y exportadores, porque una campaña con mayor producción no garantiza necesariamente un mejor margen cuando varios orígenes concentran su oferta en una misma ventana comercial \citep{ADEX2023, Velasquez2025, Zhang2025}.

La literatura reciente muestra que los métodos de inteligencia artificial y series temporales pueden contribuir a anticipar el comportamiento de los precios agrícolas. \citet{Zhang2025} aplicaron una arquitectura híbrida TCN--MLP--Attention a precios de palta y reportaron mejoras frente a modelos convencionales. \citet{Chen2023} combinaron descomposición de señales y redes LSTM para representar patrones no lineales, mientras que \citet{Xia2024} comparó ARIMA y LSTM y concluyó que el desempeño depende de la estructura, longitud y variabilidad de la serie. En conjunto, estos antecedentes sustentan la conveniencia de comparar modelos estadísticos, aprendizaje automático y aprendizaje profundo, en vez de asumir anticipadamente que una sola familia de algoritmos será superior.

En América Latina persisten brechas de digitalización y acceso a inteligencia de mercados en el sector agrícola. Los pequeños y medianos productores suelen recibir información fragmentada, tardía o condicionada por intermediarios, lo que dificulta relacionar el precio internacional con sus propios costos y rendimientos. En el caso peruano, la expansión de las exportaciones de palta ha coexistido con periodos de sobreoferta, presión sobre los precios y mayores costos logísticos. \citet{Capunay2025} describen el efecto de estas condiciones sobre la rentabilidad de las exportaciones peruanas entre 2018 y 2022, y \citet{Pantaleon2024} señalan que el crecimiento exportador no elimina la fragilidad empresarial del sector.

La Libertad constituye un ámbito relevante para esta investigación por su actividad agroexportadora y por la presencia de productores vinculados con los valles de Chao, Virú, Moche y el área de influencia de Chavimochic. El conocimiento agronómico y la experiencia de campo son indispensables, pero no permiten por sí solos procesar miles de operaciones de exportación, reconocer ciclos de precios o evaluar simultáneamente escenarios de costo, productividad y volumen. El documento técnico que sirve de base al presente proyecto identifica precisamente tres necesidades del negocio: reducir el riesgo de pérdida económica por campaña, reducir la incertidumbre mediante pronósticos de precios FOB y reducir la desinformación mediante el análisis de ciclos históricos.

La realidad problemática se concentra, por tanto, en que las decisiones de siembra, ampliación o reducción de área se adoptan frecuentemente a partir de experiencia previa y precios observados, sin una herramienta que integre datos históricos reales de exportación, costos productivos y escenarios económicos. Cuando el productor toma como referencia un precio reciente sin considerar estacionalidad, destino, volumen exportado o costos, puede sobreestimar el ingreso de la campaña. Esta limitación adquiere mayor importancia en un cultivo cuyo resultado económico depende tanto del mercado internacional como de condiciones productivas locales.

El proyecto utilizará como fuentes principales la Encuesta Nacional Agropecuaria del Instituto Nacional de Estadística e Informática (INEI) y los registros de Aduanet correspondientes a la subpartida nacional 0804400000. El documento base registra una disponibilidad aproximada de 5\,607 observaciones anuales en la fuente agropecuaria y 128\,222 registros aduaneros antes de aplicar filtros de calidad y pertinencia. Estas cantidades representan el universo disponible de registros de origen, no el tamaño definitivo del conjunto analítico. El número final se determinará después de eliminar duplicados, tratar valores ausentes, validar dominios y conservar únicamente operaciones y unidades relacionadas con el alcance de la investigación.

De Aduanet se obtendrán la fecha de operación, país de destino, valor FOB y peso neto exportado. El precio FOB unitario no se tomará directamente como un valor aislado, sino que se calculará como la relación entre el valor FOB y el peso neto válido de cada operación, para luego agregarse por semana o mes según el análisis. Del INEI se utilizarán variables de costos agrícolas, superficie, producción y cantidades equivalentes, entre ellas los campos P235\_VAL, P237\_VAL, P239, P241, P219\_CANT\_1, P219\_CANT\_2 y P219\_EQUIV\_KG, siempre que la documentación de cada encuesta confirme su correspondencia y unidad de medida. Kaggle podrá emplearse únicamente como fuente exploratoria o de contraste; no sustituirá los registros FOB peruanos.

Las variables del negocio se organizan alrededor de precio FOB por kilogramo, volumen exportado, fecha, temporada, destino, costos de producción, flete, hectáreas, productividad, margen y ganancia esperada. Las variables calculadas deberán conservar trazabilidad respecto de sus campos de origen. En particular, el margen se calculará a partir de ingresos y costos estimados, la productividad relacionará producción válida con superficie y los escenarios se construirán sobre distribuciones históricas del precio, evitando rangos arbitrarios que no estén respaldados por los datos procesados.

La solución se compone de cinco funciones analíticas principales. La primera es un módulo de cálculo de rentabilidad que integra costo de producción, precio FOB proyectado y productividad. La segunda es un simulador de escenarios optimista, esperado y pesimista que estima resultados económicos y apoya recomendaciones como sembrar, mantener, reducir área o esperar. La tercera es el modelo predictivo de precios FOB. La cuarta es un módulo interactivo para visualizar precios históricos por fecha y destino. La quinta genera reportes de ciclos, picos, caídas y comparaciones entre campañas.

La implementación preliminar compara varios enfoques en los componentes predictivos. Para precio FOB y margen se utilizaron Random Forest, HistGradientBoosting, ElasticNet y Prophet. Para la clasificación de riesgo se utilizaron Random Forest, HistGradientBoosting, regresión logística y Prophet adaptado como estimador continuo con umbral de clasificación. Además de la comparación individual, las cuatro predicciones FOB se integran mediante un ensamble de promedio simple. El módulo histórico y los reportes no entrenan modelos independientes, sino que reutilizan las observaciones y las salidas de todos los algoritmos. Esta arquitectura multimodelo combina ensambles no lineales, un modelo regularizado interpretable y un modelo temporal con regresores externos. Los modelos FOB se entrenan y evalúan con un corte cronológico común, sin mezclar observaciones futuras con el conjunto de entrenamiento.

La evaluación predictiva empleará MAE, RMSE, MAPE y $R^2$. El MAE permitirá interpretar el error promedio en la unidad del precio; el RMSE penalizará con mayor intensidad los errores grandes; y el MAPE facilitará una lectura porcentual, siempre que se controlen valores cercanos a cero. El entrenamiento preliminar se concentrará en un horizonte de seis semanas, mientras que los reportes históricos podrán presentar agregaciones mensuales.

Para el desarrollo de la solución tecnológica se consideraron tres metodologías estándar. KDD organiza el descubrimiento de conocimiento mediante selección, preprocesamiento, transformación, minería e interpretación; resulta útil para explicar la obtención de patrones, pero ofrece menor detalle para el despliegue y la retroalimentación del negocio. SEMMA estructura el trabajo en muestreo, exploración, modificación, modelado y evaluación; es adecuada para el tratamiento analítico, aunque se concentra principalmente en el proceso de modelado. CRISP-DM comprende entendimiento del negocio, entendimiento de los datos, preparación, modelado, evaluación y despliegue. Se selecciona CRISP-DM porque conecta de forma explícita las necesidades del productor con la preparación de fuentes heterogéneas, la comparación de algoritmos y la implementación de la aplicación.

La arquitectura propuesta separará la adquisición y preparación de datos, los servicios analíticos y la interfaz de usuario. El backend se desarrollará con Python y FastAPI para exponer servicios de consulta histórica, inferencia, simulación y generación de reportes. La aplicación cliente se implementará con tecnología React para entorno web adaptable o React Native si las pruebas de despliegue confirman su conveniencia móvil. La base de datos almacenará usuarios, series procesadas, ejecuciones de modelos y escenarios, procurando no conservar información personal que no sea necesaria para la investigación.

La investigación se justifica por su conveniencia, porque integra información que actualmente se consulta de forma separada; por su relevancia social, porque acerca analítica de mercado a productores con menor capacidad tecnológica; por sus implicaciones prácticas, porque traduce un pronóstico en margen y recomendaciones; por su valor teórico, porque compara cuatro familias de modelos y un ensamble bajo un mismo conjunto temporal peruano; y por su utilidad metodológica, porque documenta un flujo reproducible de extracción, validación, modelado y evaluación.

En consecuencia, el problema de investigación se formula mediante la siguiente pregunta: ¿En qué medida el desarrollo de un modelo predictivo de precios de exportación de palta Hass, integrado con módulos de rentabilidad y análisis histórico, apoyará la toma de decisiones de siembra de productores agrícolas de La Libertad, Perú?

La hipótesis general plantea que el desarrollo del modelo predictivo de precios de exportación de palta Hass, integrado con módulos de rentabilidad, escenarios y análisis histórico, mejorará el apoyo a la toma de decisiones de siembra de productores agrícolas de La Libertad. Las subhipótesis establecen que: el módulo de rentabilidad mejorará la estimación del margen esperado; el simulador reducirá la incertidumbre percibida al comparar escenarios; el modelo seleccionado alcanzará menor error que las alternativas comparadas; el análisis histórico facilitará la identificación de tendencias y estacionalidad; y los reportes de ciclos mejorarán la interpretación de periodos favorables y desfavorables.

El objetivo general es desarrollar un modelo predictivo de precios de exportación de palta Hass, integrado en una herramienta de apoyo a la toma de decisiones de siembra para productores agrícolas de La Libertad. Los objetivos específicos son: caracterizar y depurar los datos de Aduanet e INEI relacionados con precio, volumen, destino, producción y costos; construir variables temporales y económicas con trazabilidad; implementar y comparar algoritmos de regresión y pronóstico para un horizonte de seis semanas; desarrollar los módulos de rentabilidad, escenarios, análisis histórico y reportes; y evaluar la precisión del modelo y la utilidad de la herramienta con usuarios del ámbito productivo.

El estudio se limitará espacialmente a la aplicación del problema en productores de palta Hass de La Libertad y utilizará información nacional de exportación y producción que permita contextualizar dicho ámbito. Temporalmente, procesará los registros disponibles y verificables comprendidos entre 2018 y el último periodo completo accesible al momento de cerrar la base. No se afirmará que el sistema predice el precio recibido individualmente por cada productor, pues el valor FOB agregado puede diferir del precio en chacra por contratos, calidad, calibre, intermediación y logística. Tampoco se presentarán como resultados ejecutados las evaluaciones de usuarios o métricas de modelos que todavía correspondan a actividades futuras del proyecto.
